\section{The true section occupation under a local collision constraint}
\label{sec:v37-occupation-local}
We return to the original circular family and the actual section
$Y_R^*$ of \eqref{eq:enlarged-section}. The occupation of this section
is not the constant collision step. Nevertheless its geometry permits
a bounded multiplier at one collision, and hence a perturbation at
occupation frequency of order $m^{-1/2}$. This gives a joint local and
central limit without changing the section or the physical record.

Put $c=c_*$, $\eta_R=\one_{Y_R^*}$, and define
\begin{equation}\label{eq:v37-occupation-record}
 A_{m,R}=\sum_{j=0}^{m-1}\eta_R\circ T_R^j,\quad
 E_{m,R}=A_{m,R}-mc,\quad
 X_R=(\kappa_{R,1},\kappa_{R,2},\tau_R-\bar\tau_R,\eta_R-c).
\end{equation}
The sum includes the initial collision and excludes collision $m$.
The symbol $A_{m,R}$ is an occupation count, not a return time.

\subsection{An exact, unsmoothed occupation multiplier}
\begin{lemma}[The actual rectangular section is a multiplier]
\label{lem:v37-section-multiplier}
Multiplication by $\eta_R$ is bounded on every local strong collision
space, uniformly in $R$. Consequently, for $v\in\mathbb C$,
\begin{equation}\label{eq:v37-occupation-exponential}
 M_{e^{iv(\eta_R-c)}}=e^{-ivc}\{I+(e^{iv}-1)M_{\eta_R}\}
\end{equation}
is an entire strong-space family with uniform bounds on compact sets.
No smoothing of $\eta_R$ occurs in this identity.
\end{lemma}
\begin{proof}
The section is the union of five coordinate rectangles, with four of
them small squares, in $(\alpha,p)$ coordinates. Their horizontal
boundaries remain horizontal after $p=\sin\varphi$, and their vertical
boundaries remain vertical after the local constant arclength
reparametrization. All lie in a common compact region away from grazing.
Split the angular circle at a fixed meridian and partition by all
rectangle endpoints in both coordinates. This gives a uniformly finite
partition into open, simply connected cells on which $\eta_R$ is
constant. It is a partition of the collision cylinder modulo null
boundaries, not a partition of a long orbit.

Admissible stable curves are graphs whose slopes are negative and
bounded away from zero and infinity in these local coordinates.
They therefore cross any horizontal or vertical partition side at
most once. Their intersection with a rectangular cell is an interval.
Their length within distance $\epsilon$ of one such side is bounded
by $C\epsilon$, uniformly in the location of the side. Summing over
the fixed finite number of sides gives the partition conditions
(1)--(2) of \cite[Lemma 3.3]{DPZ}, uniformly on the finite parameter
cover. The constant functions on the cells meet every strict
piecewise-H\"older exponent requirement. Part (b) of that lemma gives
$\|M_{\eta_R}\|\le C$. Formula
\eqref{eq:v37-occupation-exponential} follows from $\eta_R^2=\eta_R$
pointwise and proves the remaining assertions.
\end{proof}

Let $Q_{R,z}$ be the exact displacement--roof operator of
Section~\ref{sec:compact-collision-spectrum}, and set
\begin{equation}\label{eq:v37-joint-operator}
 \mathcal Q_{R,z,v}=Q_{R,z}M_{e^{iv(\eta_R-c)}}.
\end{equation}
The occupation multiplier is on the source side of each collision.
If $a\nu$ belongs to the strong space and $M_d$ is bounded there,
then for bounded physical representatives $a,d$,
\begin{align}\label{eq:v37-exact-joint-pairing}
 \ell\bigl(M_d\mathcal Q_{R,z,v}^{m}(a\nu)\bigr)
  =\int a(x)d(T_R^mx)
   e^{iz\cdot(K^{\rm c}_{m,R},S_{m,R}-m\bar\tau_R)+ivE_{m,R}}d\nu(x).
\end{align}
The equality follows by the transfer convention at each step. Products
of the finitely many bounded multipliers are understood as physical
weighted compositions first and then by continuity. This also follows
from the module property of the piecewise multiplier and its dual
pairing with smooth tests. In particular it applies to $a,d$ smooth,
to $a=d=\eta_R$, and to their products with smooth functions.
The distinction between the source occupation and the image-side
preceding roof in \eqref{eq:v37-joint-operator} fixes the endpoint count.

\begin{proposition}[Joint expansion and a small-occupation-frequency gap]
\label{prop:v37-joint-small-frequency}
On a uniform complex neighborhood of zero, $\mathcal Q_{R,z,v}$ has
a simple eigenvalue $\lambda_R(z,v)$ and a rank-one projection,
with exponentially decaying complementary powers. Writing
$\Theta=(z,v)\in\mathbb C^4$,
\begin{equation}\label{eq:v37-joint-expansion}
 \log\lambda_R(\Theta)=-\tfrac12\Theta^{\mathsf T}\Omega_R\Theta
                         +O(|\Theta|^3),
\end{equation}
with constants uniform in $R$. The matrices $\Omega_R$ are continuous
and uniformly positive definite. For each $B<\infty$ and $\epsilon>0$
there are $v_{B,\epsilon}>0$, $C_{B,\epsilon}<\infty$ and
$\rho_{B,\epsilon}<1$ such that
\begin{equation}\label{eq:v37-compact-small-occupation}
 \|\mathcal Q_{R,\omega,v}^m\|
       \le C_{B,\epsilon}\rho_{B,\epsilon}^m
 \quad (|b|\le B,\ |\omega|\ge\epsilon,\ |v|<v_{B,\epsilon}),
\end{equation}
where $\omega=(u,b)\in\T^2\times\R$ is real and $v$ may be complex.
The occupation radius is chosen after $B$ and $\epsilon$.
\end{proposition}
\begin{proof}
Lemma~\ref{lem:v37-section-multiplier} and the exact preceding-flight
multiplier give strong analyticity with common bounds on a fixed
complex ball. Perturb the uniform unweighted splitting by its usual
resolvent contours. At zero, the first eigenvalue derivative vanishes
because every component of $X_R$ is centered. The second derivative
of the logarithm is minus the asymptotic covariance of the physical
sum in \eqref{eq:v37-exact-joint-pairing}. One can identify it directly:
differentiate that finite-time pairing twice, divide by $m$, and use
the analytic splitting. The amplitude derivatives divided by $m$ tend
to zero and the complementary derivatives decay exponentially, by
Cauchy's estimate on a smaller fixed ball.

There is an exact collision identity $X_R=L_Rh_R$, where
\begin{equation}\label{eq:v37-covariance-change}
 L_R=\begin{pmatrix}
 1&0&0&0\\0&1&0&0\\0&0&-\bar\tau_R&1\\0&0&-c&0
 \end{pmatrix},\qquad
 \Omega_R=L_R\Gamma_RL_R^{\mathsf T}=cL_RD_RL_R^{\mathsf T}.
\end{equation}
Here $h_R$ is the bounded exact compensation, not an induced observable.
The covariance identification and convergence follow from
Theorem~\ref{thm:collision-covariance}. Since $\det L_R=c>0$,
Theorem~\ref{thm:joint-uniform-nondegeneracy} implies uniform positivity
and continuity of $\Omega_R$. Uniform analyticity gives the cubic
remainder in \eqref{eq:v37-joint-expansion}.

For the last assertion first fix $B,\epsilon$. Theorem~\ref{thm:compact-collision-spectrum}
bounds $\|Q_{R,\omega}^m\|\le C_0\rho_0^m$ on this compact set.
Choose $\rho_0<\rho_1<1$. The series for
$(w-Q_{R,\omega})^{-1}$ is uniformly bounded on $|w|=\rho_1$
and on its exterior up to a common operator-norm bound. By
\eqref{eq:v37-occupation-exponential},
$\|\mathcal Q_{R,\omega,v}-Q_{R,\omega}\|\le C_B|v|$ for small $v$.
Choose $v_{B,\epsilon}$ so the resolvent Neumann series converges
uniformly on that exterior. Cauchy's power formula on $|w|=\rho_1$
proves \eqref{eq:v37-compact-small-occupation}. No weighted
Lasota--Yorke estimate for an entire occupation-frequency torus has
been used.
\end{proof}

Write the positive covariance and its Schur complement as
\begin{equation}\label{eq:v37-schur-data}
 \Omega_R=\begin{pmatrix}\Sigma_R&b_R\\b_R^{\mathsf T}&d_R\end{pmatrix},
 \qquad \beta_R=b_R^{\mathsf T}\Sigma_R^{-1},\qquad
 \sigma_{A,R}^2=d_R-b_R^{\mathsf T}\Sigma_R^{-1}b_R>0.
\end{equation}
All these matrices and scalars are continuous; their relevant positive
bounds are uniform. For $z\in\R^3$ let $\mathcal N_{R,z}$ be the real
normal law with mean $\beta_Rz$ and variance $\sigma_{A,R}^2$.

\subsection{A local roof interval and a central occupation variable}
An endpoint family below is called multiplier-bounded if $a_R,d_R^{\rm end}$
are bounded physical functions, $\|a_R\nu\|_{\mathcal B}$ and
$\|M_{d_R^{\rm end}}\|$ are bounded on the local spaces, and their
supremum norms are uniformly bounded. Smooth bounded families with
common $C^2$ bounds and $a_R=d_R^{\rm end}=\eta_R$ are examples.
The endpoint notation $d_R^{\rm end}$ is distinct from the covariance
entry $d_R$ in \eqref{eq:v37-schur-data}.

\begin{theorem}[Mixed local and central limit for actual occupation]
\label{thm:v37-occupation-local-central}
Let $a_R,d_R^{\rm end}$ be nonnegative multiplier-bounded endpoint
families, and put $\alpha_R=\nu(a_R)\nu(d_R^{\rm end})$. Fix a bounded
interval $J\subset\R$ of positive length and a bounded continuous
function $F:\R\to\mathbb C$. Uniformly for $R\in I$ and central
targets with $|z_{m,R}(k,t)|\le M$, where
$z_{m,R}(k,t)=(k/\sqrt m,(t-m\bar\tau_R)/\sqrt m)$,
\begin{align}
 &m^{3/2}\int a_R(x)d_R^{\rm end}(T_R^mx)
  \one_{\{K^{\rm c}_{m,R}=k,\ S_{m,R}-t\in J\}}
                    F(E_{m,R}/\sqrt m)\,d\nu(x)\notag\\
 &\hspace{12mm}
 -\alpha_R|J|g_{\Sigma_R}(z_{m,R}(k,t))
                    \int F\,d\mathcal N_{R,z_{m,R}(k,t)}
                   \longrightarrow0.             \label{eq:v37-mixed-local-central}
\end{align}
The assertion also holds for the indicator of any fixed interval in
the occupation variable, including a half-line. In particular the
conditional occupation law is $\mathcal N_{R,z}$ in the limit whenever
$\alpha_R$ is uniformly positive. The conditioned event in this
assertion is unchanged.
\end{theorem}
\begin{proof}
First replace $\one_J(S_m-t)$ by an integrable time test $q(S_m-t)$
with integrable Fourier transform supported in $[-B,B]$, and take
$F(y)=e^{i\theta y}$. The exact Fourier expression contains
$\widehat q(-b)$, torus orthogonality in the two displacement
coordinates, and the operator
$\mathcal Q_{R,(u,b),\theta/\sqrt m}^m$. Fix $B$ and a compact set of
$\theta$ before taking $m\to\infty$.
Outside a fixed neighborhood of $\omega=0$, use
\eqref{eq:v37-compact-small-occupation}. Inside it, use
\eqref{eq:v37-joint-expansion}, the principal amplitude
$\alpha_R+O(|\omega|+|\theta|/\sqrt m)$, and put $w=\sqrt m\,\omega$.
Uniform positivity gives an integrable Gaussian bound in $w$ for
bounded $\theta$; the complementary power is exponentially small.
Dominated convergence yields the partial inverse
\begin{align}\label{eq:v37-partial-gaussian-inverse}
 &(2\pi)^{-3}\int_{\R^3}e^{-iw\cdot z}
        e^{-\frac12(w,\theta)^{\mathsf T}\Omega_R(w,\theta)}\,dw\notag\\
 &\hspace{12mm}=g_{\Sigma_R}(z)
       \exp\{i\theta\beta_Rz-\tfrac12\theta^2\sigma_{A,R}^2\}.
\end{align}
Multiplying by $\alpha_R\int q$ proves the band-limited assertion,
uniformly in central targets. The sign of the conditional mean in
\eqref{eq:v37-partial-gaussian-inverse} follows from the inverse phase
$e^{-iw\cdot z}$.

Next remove the time test. Let $q_B^-\le\one_J\le q_B^+$ be the
positive-envelope pair used in Lemma~\ref{lem:collision-interval-envelopes}.
Their Fourier supports are fixed and
$\int(q_B^+-q_B^-)=O(B^{-1})$. Although the occupation phase is
complex, its modulus is one. Thus the absolute replacement error
using $q_B^-$ is bounded by the same nonnegative endpoint integral
with time test $q_B^+-q_B^-$ and with occupation frequency zero.
The already proved band-limited assertion at zero bounds the limsup
of the scaled error by $C/B$. Take $m\to\infty$ for each fixed $B$
and then $B\to\infty$. This proves characteristic-function convergence
for the exact interval $J$, also at $\theta=0$. A singleton endpoint
of $J$ is handled by majorants of intervals of arbitrarily small
length, so the choice of open or closed endpoints has no effect.

The positive finite measures in the occupation variable now have
converging masses and Fourier transforms continuous at zero, with
the Gaussian limits in \eqref{eq:v37-partial-gaussian-inverse}.
The continuity theorem for finite positive measures gives weak
convergence. If a subsequence of masses tends to zero, boundedness of
$F$ handles it directly. Otherwise normalize by the masses and apply
the probability version. This proves the assertion for bounded
continuous $F$ and for Gaussian continuity sets. Uniformity follows
by contradiction: along a violating sequence take convergent
subsequences of $R$, $z$, and the bounded amplitudes $\alpha_R$;
the Fourier proof and covariance continuity identify every limit as
the displayed Gaussian finite measure. For a uniformly positive
$\alpha_R$, division by its exact local mass proves the conditional
claim. No event replacement and no growing-band spectral estimate
enters the argument.
\end{proof}

\begin{remark}[What is and is not estimated]
The occupation inserted in \eqref{eq:v37-joint-operator} is the exact
$\eta_R$, not a surrogate chosen after a smoothing limit. The estimate
uses the entire fixed displacement--roof band and an occupation
frequency shrinking to zero. It does not estimate fixed nonzero
occupation frequencies, or a pointwise roof density. These two
additional inversions are not part of
\eqref{eq:v37-mixed-local-central}.
\end{remark}
